% -----------------------------------------------
% Template for ISMIR Papers
% 2026 version, based on previous ISMIR templates

% Requirements :
% * 6+n page length maximum
% * 10MB maximum file size
% * Copyright note must appear in the bottom left corner of first page
% * Clearer statement about citing own work in anonymized submission
% (see conference website for additional details)
% -----------------------------------------------

\documentclass{article}
\usepackage[T1]{fontenc}
\usepackage[utf8]{inputenc}
\usepackage{ismir}
\usepackage{amsmath,amssymb,cite,url}
\usepackage{graphicx}
\usepackage{color}
\graphicspath{{../Draft-legacy/v2-0213-2026/V2-latex/figures/}}

\title{When Does Model Capacity Matter? Evaluating Pretrained Audio Representations for Music Retrieval in Small-Scale Classical Archives}

\oneauthor
  {Keita Katsumi}
  {California State Polytechnic University, Pomona\\\texttt{kkatsumi@cpp.edu}}

\def\authorname{K. Katsumi}

\setlength{\parskip}{0pt}
\setlength{\parindent}{1em}

\sloppy % please retain sloppy command for improved formatting

\begin{document}

\maketitle

\begin{abstract}
This study investigates the effect of embedding-model capacity on content-based music similarity retrieval in a small-scale classical music archive. While larger models are often assumed to yield superior representations, their effectiveness in limited, single-domain settings remains unclear. We evaluate CNN-based and Transformer-based embedding models across multiple capacity tiers using two proxy tasks: composer-based retrieval and musical-character retrieval. Retrieval performance is assessed using ranking-based metrics (NDCG, Hit@K) and set-based metrics (precision, recall, and F1). Results show that increasing model capacity does not consistently improve retrieval performance. While higher-capacity models may provide limited gains for structured tasks, these effects are not consistent across model families, and scaling effects are weak or absent for abstract semantic retrieval. Furthermore, we observe that evaluation outcomes are strongly influenced by the relationship between the number of relevant items and the evaluation cutoff. In particular, when a large number of relevant items exist per query, set-based metrics such as Recall and F1 provide limited signal at small cutoff values, while ranking-based metrics such as NDCG and Hit@K remain more stable. These findings suggest that retrieval performance in small-scale archival settings is governed more by task structure and evaluation design than by model capacity alone.
\end{abstract}

% Section 1 Introduction
\section{Introduction}\label{sec:introduction}
Modern music streaming platforms rely heavily on machine learning to support search, recommendation, and discovery. Much of the prior work in this area, as well as many industrial systems, is built around large-scale collaborative filtering and deep models trained on extensive user interaction logs. While effective for mainstream catalogs, these approaches are ill-suited to small, curated archives that lack dense user data and whose mission emphasizes exploration over pure popularity. Classical music archives such as the iPalpiti collection exemplify this setting: they contain high-quality, stylistically coherent recordings but limited usage data, and they require tools that can surface musically meaningful relationships among works, performances, and artists.

In such environments, content-based recommendation driven by audio embeddings becomes central.
Recent advances in pretrained audio models---including lightweight convolutional networks and large transformer-based architectures---enable the extraction of rich musical representations from raw audio. However, most evaluations of these models are conducted on large-scale, heterogeneous datasets, where increased model capacity often correlates with improved benchmark performance.
It remains unclear whether this trend transfers to small, single-domain archives such as classical
music collections. In small-scale settings, limited data diversity, coarse proxy definitions, and the presence of large relevant item sets under multi-label similarity may reduce the practical benefit of increased model complexity. This raises a fundamental design question for content-based recommendation systems in archival contexts: does increasing embedding-model capacity meaningfully improve similarity retrieval performance, or does it introduce diminishing returns under constrained conditions? This study further investigates whether the effect of model capacity is conditioned by the structure of the retrieval task itself.

This reseach addresses this question by focusing exclusively on the backend: the stage that maps audio recordings to embeddings and supports similarity search, independent of end-to-end user modeling. Specifically,the study adopts a hypothesis-driven, iterative experimental design.
Rather than assuming that larger models necessarily yield better retrieval performance, we begin with the hypothesis that increased embedding-model capacity may not improve similarity retrieval in small-scale, single-domain archives.

To evaluate this hypothesis, retrieval performance is systematically assessed across a spectrum of
model capacities within two representative architectural families (CNN-based and Transformer-based
models). Proxy retrieval tasks are formulated to approximate musically meaningful similarity, including both a sanity-check retrieval task (composer-based) and a musical-character-based proxy (e.g., contrasting musical characteristics such as \emph{energetic} vs.\ \emph{calm}) Ranking quality is evaluated using standard information-retrieval metrics, including both ranking-based metrics (e.g., NDCG, Hit@K) and set-based metrics (e.g., Precision@K, Recall@K). This distinction allows us to examine how evaluation design interacts with multi-label relevance and the relationship between the number of relevant items and the evaluation cutoff in small-scale retrieval settings. Experimental conditions are progressively expanded to examine whether higher-capacity models provide consistent improvements, exhibit diminishing returns, or fail to deliver measurable gains under constrained archival settings.This research is guided by two research questions:

\smallskip
\noindent\textbf{RQ1:} In a small-scale classical music archive (iPalpiti), how does embedding-model capacity influence content-based music similarity retrieval performance?

\noindent\textbf{RQ2:} Under what conditions, if any, do higher-capacity embedding models provide
meaningful gains over lighter architectures?

\noindent\textbf{Hypothesis:} In a small-scale, single-domain classical music dataset, increasing
embedding-model capacity does not necessarily improve music-similarity retrieval performance.

\smallskip
This work makes three main contributions. First, it provides a systematic empirical investigation of the relationship between embedding-model capacity and content-based similarity retrieval performance in a small-scale, single-domain classical archive. Second, it establishes musically grounded evaluation criteria (via structured proxy retrieval tasks) that enable the identification of conditions under which retrieval improvements are meaningful. Third, through a hypothesis-driven experimental design, it analyzes how embedding-model capacity behaves under constrained archival conditions, clarifying the practical limits of model complexity in small, single-domain
recommendation settings.

% Section 2 Related Work
\section{Related Work}\label{sec:related-work}

Content-based recommendation is especially important in music collections with sparse or absent user interaction data, where embeddings derived from the audio signal support similarity search,
retrieval, and downstream recommendation. Recent work has shown the effectiveness of pretrained audio representations in music information retrieval while also identifying a gap in their systematic evaluation within recommender systems \cite{tamm2024}. That gap is particularly relevant for small, domain-specific archives such as classical collections, whose stylistic coherence, limited scale, and rich structural relationships raise transferability questions for models developed on large, heterogeneous datasets.

Prior work has explored a broad range of audio representation architectures, including convolutional neural networks (CNNs), recurrent neural networks (RNNs), transformer-based models, and hybrid architectures \cite{zaman2023}. Historically, CNN and hybrid CNN–RNN models have dominated supervised music tasks \cite{jiang2024}, while more recent transformer-based approaches use self-attention to model longer-range structure \cite{zaman2023}. In this reseach, however, the experimental design focuses specifically on pure CNN-based and pure transformer-based model families. CNN-based models remain attractive because their inductive biases support efficient learning of local spectral structure, and compact CNN architectures have shown strong performance in music detection and related tasks \cite{kong2020panns}. Hybrid architectures have also been explored, including CNN–RNN models for temporal modeling \cite{lin2024} and CNN–Transformer models for combining local and long-range representations \cite{pourmoazemi2024}, but they are treated here as background rather than as primary comparison targets.

Transformer-based models provide a contrasting representation regime: AST applies self-attention
directly to spectrograms \cite{gong2021ast}. Although these models can deliver strong downstream
performance, their advantages are typically established on large and diverse datasets, leaving open the question of how well such capacity transfers to small, domain-constrained classical archives. Evaluation strategy is equally important. Classification benchmarks remain common for audio embeddings \cite{gong2021ast, zaman2023}, but strong classification accuracy does not guarantee well-structured neighborhoods for similarity retrieval \cite{tamm2024}. For retrieval-oriented settings, ranking metrics such as NDCG as well as set-based metrics such as Precision@K and Recall@K, are commonly used to evaluate ranking quality \cite{krichene2020}. However, these metrics can behave differently under multi-label relevance conditions: set-based metrics may provide limited or unstable signals when the number of relevant items per query is large relative to the evaluation cutoff, whereas ranking-based metrics tend to remain more robust. This limitation is consistent with  prior observations in MIR research, where datasets are often small or domain-specific, posing challenges for generalizable evaluation \cite{urbano2013}.At the same time, these metrics are sensitive to evaluation design: limited candidate pools can reduce discriminative power, increase ties, and even alter comparative conclusions, particularly when many relevant items exist per query but only a small top-K set is evaluated \cite{krichene2020, canamares2020}; this is especially important in small archival collections. Diversity and long-tail considerations also remain relevant background concerns in music recommendation, particularly in settings where exploration is valued\cite{porcaro2022}.

Broader recommender-system work addresses personalization through user modeling, sequential
interaction modeling, and contextual signals \cite{lin2025, abbattista2024, schedl2021}, but those
concerns are outside the scope of this research, which isolates backend embedding quality from user modeling and large-scale deployment constraints. Overall, prior work has established that larger and more complex models often improve performance in data-rich settings, yet it remains unclear whether embedding-model capacity provides comparable gains in small, single-domain classical archives. This research addresses that gap by examining how capacity interacts with retrieval-oriented evaluation under constrained archival conditions \cite{tamm2024, krichene2020, canamares2020}.

% Section 3 Method
\section{Method}\label{sec:method}\label{ch:problem}\label{ch:methodology}
This section formalizes the retrieval setting and experimental framework used to evaluate the effect of embedding-model capacity on similarity-based retrieval performance.

\subsection{Problem Formulation}

Let $D = \{x_1, x_2, \dots, x_N\}$ denote a collection of $N$ classical music recordings drawn from
the iPalpiti archive. Each item $x_i$ represents a complete audio track.

An embedding model is defined as a function:
\[
f_\theta : \mathcal{X} \rightarrow \mathbb{R}^d
\]
where $\theta$ denotes the model parameters and $d$ the embedding dimensionality. The embedding of
item $x_i$ is:
\[
z_i = f_\theta(x_i)
\]

Similarity between two items is computed using a similarity function $s(\cdot,\cdot)$, such as cosine
similarity:
\[
s(z_q, z_i) = \frac{z_q \cdot z_i}{\|z_q\|\|z_i\|}
\]
Higher similarity values indicate greater embedding-space proximity.

Given a query item $x_q \in D$, all items in the collection except $x_q$ are ranked in descending
order of similarity $s(z_q, z_i)$. Retrieval quality is evaluated with respect to a relevance
function $r_q(x_i) \in \{0,1\}$ defined by proxy tasks, and ranking quality is measured using both ranking-based metrics (e.g., NDCG@K, Hit@K) and set-based metrics (e.g., Precision@K, Recall@K, and F1@K), enabling analysis of how different evaluation criteria behave under multi-label relevance and varying relationships between the number of relevant items and the evaluation cutoff.

The primary independent variable is embedding-model capacity, operationalized mainly through
parameter count ($|\theta|$), with architectural class treated as a separate categorical factor.
Capacity effects are examined within architectural families via small, base, and large variants under a controlled retrieval setting.

Relevance is defined through two metadata-based proxy tasks. Let $g(x)$ denote a metadata labeling
function.

\textbf{Sanity Proxy.} Relevance is defined by shared composer identity:
\[
r_q(x_i) = 1 \quad \text{if} \quad g_{\text{composer}}(x_q) = g_{\text{composer}}(x_i)
\]

\textbf{Primary Musical Proxy.} Relevance is defined by shared character labels:
\[
r_q(x_i) = 1 \quad \text{if} \quad g_{\text{character}}(x_q) \cap g_{\text{character}}(x_i) \neq \emptyset
\]

For experimental tractability, relevance is treated as binary by projecting multi-label annotations into a match-based criterion:
\[
r_q(x_i) \in \{0,1\}
\]
where $r_q(x_i)=1$ if at least one label is shared between $x_q$ and $x_i$, and $r_q(x_i)=0$ otherwise. This formulation simplifies evaluation while preserving the underlying multi-label structure of the data. This study assumes that such metadata-derived labels provide a controlled, though coarse, approximation of musical similarity.

\subsection{Retrieval Pipeline}
The backend evaluation pipeline consists of four stages: audio preprocessing, embedding extraction, similarity computation, and ranking. All components other than embedding extraction are held constant to isolate the effect of embedding capacity.

The dataset is the iPalpiti classical music archive, comprising 203 tracks and 24.9 hours of
audio. Ground truth is derived from editorial metadata, including composer identity and character
labels used in the proxy tasks.

Each complete track serves as the retrieval unit, and embeddings are computed directly from
full-length audio inputs without segmentation. Cosine similarity is used for ranking across all
models. This controlled design ensures that performance differences are attributable primarily to
the embedding representations.

\subsection{Model Setup}
Models are selected to represent discrete capacity tiers within two architectural families: CNN-based and Transformer-based models. Within each family, small, medium, and large variants enable controlled scaling comparisons, while cross-family comparisons are interpreted descriptively because architecture, parameter count, and pretraining objectives may co-vary.

CNN-based models are drawn from a unified convolutional lineage in which scaling is achieved
primarily through increased depth under comparable training objectives. Specifically, we use Cnn6 (small), Cnn10 (medium), and Cnn14 (large), following the PANNs framework proposed by Kong et al. \cite{kong2020panns}, to represent controlled capacity scaling within the CNN family.

Transformer-based models are selected from a self-supervised audio representation framework
pretrained on large-scale music data. We use MERT-95M (medium) and MERT-330M (large) 
to represent different capacity levels within the transformer family, following the MERT framework \cite{li2023mert}.

Note that parameter counts are not strictly aligned across architectural families; therefore, cross-family comparisons are interpreted descriptively rather than as direct capacity-matched evaluations. All selected models provide pretrained checkpoints and support standardized embedding extraction.

\subsubsection{Model Capacity Definition}
In this study, model capacity is defined based on architectural depth and model scale within each family.

For CNN-based models, capacity is controlled by the number of convolutional layers following the PANNs architecture\cite{kong2020panns}:
\begin{itemize}
  \item CNN-Small (Cnn6, 6 layers)
  \item CNN-Medium (Cnn10, 10 layers)
  \item CNN-Large (Cnn14, 14 layers)
\end{itemize}

For Transformer-based models, capacity is defined by pretrained parameter scale:
\begin{itemize}
  \item Transformer-Medium (MERT-95M)
  \item Transformer-Large (MERT-330M)
\end{itemize}

These variants differ primarily in network depth, which determines model capacity within the CNN family. Although parameter counts could provide an alternative definition of capacity, we define capacity within each architectural family based on standard configurations (e.g., depth for CNNs and pretrained model scale for Transformers), as parameter counts are not directly comparable across architectures.

\subsection{Label Construction}
Relevance labels for the proxy tasks are derived from metadata and a semi-automatic annotation pipeline.

For the sanity proxy, composer labels are normalized to a canonical form to ensure consistent
matching across recordings. This normalization is performed during pseudo-label generation and
stored as a stable evaluation field.

For the primary musical proxy, character labels are constructed using a semi-automatic pipeline.
Initial label candidates are defined based on the arousal–valence framework \cite{eerola2011}, which represents emotion in a continuous dimensional space. To support this process, pretrained music emotion recognition models are used as auxiliary tools for generating candidate labels, grounded in recent advances in deep learning-based music emotion recognition \cite{jiang2024}.

Ambiguous cases are manually reviewed to ensure consistency across the dataset.
The resulting pseudo-labels are fixed prior to evaluation and reused across all experiments.

\subsection{Evaluation Protocol}
Each track in the archive is used as a query in a leave-one-out retrieval setting, where the query
item itself is excluded from the candidate pool. Relevance is determined by metadata matching as
defined above, and performance metrics are averaged across all query instances.

Because relevance labels are binary and the dataset is relatively small, ranking ties may occur when multiple relevant items share identical similarity scores. To ensure deterministic evaluation, a consistent tie-breaking strategy (e.g., lexicographic ordering by track identifier) is applied for determinism, although we acknowledge that tie-breaking may influence ranking-based metrics in small datasets.

Performance is primarily assessed using ranking-based metrics (NDCG@5, Hit@5) alongside set-based metrics (Precision@5, Recall@5, F1@5), with interpretation focusing on consistent trends across metrics, particularly under conditions where the number of relevant items is large relative to the evaluation cutoff.

Furthermore, small numerical differences in mean scores do not necessarily correspond 
to meaningful differences in user experience. Prior work emphasizes the importance 
of distinguishing between statistical and practical significance in IR evaluation 
\cite{urbano2013}.

% Section 4 Results
\section{Results}\label{sec:results}\label{ch:results}
\subsection{Sanity Proxy Results}

The Sanity Proxy task evaluates the ability of embedding models to retrieve tracks sharing the same composer. This task serves as a baseline for verifying that models capture fundamental compositional structures.

Table~\ref{tab:sanity_results} presents retrieval performance grouped by architectural family and ordered by capacity tier within each family.

\begin{table}[!htbp]
\centering
\caption{Sanity Proxy Results (Composer Retrieval). Metrics reported are corpus-level mean values at rank 5.}
\label{tab:sanity_results}
\resizebox{\columnwidth}{!}{%
\begin{tabular}{lccccc}
\hline
\textbf{Model} & \textbf{NDCG@5} & \textbf{Hit@5} & \textbf{Precision@5} & \textbf{Recall@5} & \textbf{F1@5} \\
\hline
\multicolumn{6}{c}{\textbf{CNN Family}} \\
\hline
CNN-Small & 0.548 & 0.640 & 0.266 & 0.195 & 0.188 \\
CNN-Medium & 0.545 & 0.640 & 0.303 & 0.215 & 0.214 \\
CNN-Large & 0.585 & 0.660 & 0.321 & 0.228 & 0.229 \\
\hline
\multicolumn{6}{c}{\textbf{Transformer Family}} \\
\hline
Transformer-Medium & 0.642 & 0.709 & 0.361 & 0.263 & 0.265 \\
Transformer-Large & 0.588 & 0.665 & 0.304 & 0.233 & 0.224 \\
\hline
\end{tabular}
}
\end{table}

The Sanity Proxy results show that increased capacity can improve retrieval when relevance is defined by a relatively structured and objective criterion. 

Within the CNN family, NDCG@5 drops slightly from the small to the medium model before showing notable improvement with the large model, suggesting that sufficient capacity scaling is required to see benefits. Conversely, within the Transformer family, performance degrades substantially from the medium to the large model across all metrics. 

\subsection{Primary Musical Proxy Results}
The Primary Musical Proxy task evaluates retrieval based on abstract musical character labels (e.g., "Energetic" vs. "Calm"). This task represents the core retrieval challenge for the archive.

Table~\ref{tab:musical_results} summarizes the performance across all models.

\begin{table}[!htbp]
\centering
\caption{Primary Musical Proxy Results (Character Retrieval). Metrics reported are corpus-level mean values at rank 5.}
\label{tab:musical_results}
\resizebox{\columnwidth}{!}{%
\begin{tabular}{lccccc}
\hline
\textbf{Model} & \textbf{NDCG@5} & \textbf{Hit@5} & \textbf{Precision@5} & \textbf{Recall@5} & \textbf{F1@5} \\
\hline
\multicolumn{6}{c}{\textbf{CNN Family}} \\
\hline
CNN-Small & 0.653 & 0.783 & 0.536 & 0.039 & 0.071 \\
CNN-Medium & 0.656 & 0.768 & 0.528 & 0.038 & 0.070 \\
CNN-Large & 0.642 & 0.749 & 0.528 & 0.038 & 0.070 \\
\hline
\multicolumn{6}{c}{\textbf{Transformer Family}} \\
\hline
Transformer-Medium & 0.632 & 0.778 & 0.447 & 0.033 & 0.060 \\
Transformer-Large & 0.631 & 0.764 & 0.485 & 0.036 & 0.066 \\
\hline
\end{tabular}
}
\end{table}

In contrast to the Sanity Proxy, the Musical Proxy results do not show a clear benefit from increased capacity. While Recall@5 and F1@5 values are extremely low across all models (approximately 0.03--0.04), Hit@5 remains relatively high (around 0.75), and NDCG@5 remains stable. This indicates that relevant or partially relevant items are often retrieved within the top-ranked results, but do not fully match the strict set-based relevance criteria.

This behavior arises from the structure of the task: each query is associated with a relatively large set of relevant items due to multi-label tag assignments, while evaluation is performed at a small cutoff (K=5). As a result, even when relevant or partially relevant items are retrieved near the top ranks, only a small fraction of the total relevant set is captured, leading to low Recall and F1 scores.

Despite this evaluation characteristic, the overall trend remains consistent: increasing model capacity does not yield meaningful improvements across models. Both CNN and Transformer families exhibit similar performance ranges, supporting the hypothesis that capacity effects are strongly task-dependent in this setting.

\subsection{Capacity--Performance Trends Within Architectural Families}
This section analyzes scaling behavior strictly within architectural families to isolate the effect of parameter count from architectural inductive bias. Taken together, the within-family scaling patterns described below provide the primary empirical test of the condition-dependent capacity hypothesis.

\subsubsection{CNN Scaling Behavior}
Within the CNN family, capacity scaling is task-dependent. On the Sanity Proxy, performance remains comparable between the small and medium tiers and improves at the large tier, indicating that modest increases in capacity can help for structured retrieval but that further scaling yields limited gains. On the Musical Proxy, however, the pattern is non-monotonic.

\subsubsection{Transformer Scaling Behavior}
Within the Transformer family, the effect of scaling also depends on the proxy task. On the Sanity Proxy, the medium transformer outperforms the large model, indicating that increased capacity does not necessarily lead to improved performance even when the task is aligned with structured musical metadata. On the Musical Proxy, however, the medium and large transformer models perform comparably, indicating that increased capacity does not translate into improved abstract semantic retrieval. This contrast reinforces the view that transformer scaling benefits are conditional rather than universal in the present setting.

\subsection{Cross-Family Descriptive Comparison}
Comparing performance across families reveals that scaling effects depend more on task structure
than on raw model size alone. On the Sanity Proxy, both families show some sensitivity to capacity scaling, although the effect is not strictly monotonic. In particular, improvements are observed within the CNN family, while the Transformer family does not exhibit consistent gains with increased capacity. On the Musical Proxy, however, neither family exhibits consistent improvement, and performance differences remain small.

These comparisons are descriptive and do not imply causal superiority of one architecture over
another, since architecture, parameter count, and pretraining regime co-vary across families.
The narrow performance range across models suggests limited representational separability under
the current retrieval setting, regardless of capacity tier. Given the small dataset size, absolute differences should be interpreted cautiously, with emphasis placed on consistent trends rather than isolated values.

\subsection{Summary of Empirical Findings}
The experimental results yield three main observations. First, increasing model capacity does not guarantee monotonic performance improvement. While some gains are observed in the Sanity Proxy, these gains either saturate or remain limited to specific settings; on the Musical Proxy, scaling effects are weak or absent. Second, scaling behavior differs clearly by task. Structured composer retrieval benefits more from increased capacity, whereas abstract character-based retrieval does not show consistent improvements. Third, performance values cluster within a relatively narrow range across models, suggesting limited separability under the current retrieval setup. Taken together, these findings indicate that the effect of embedding-model capacity is strongly task-dependent. While increased capacity can improve retrieval when relevance is defined by structured and well-constrained attributes, it does not reliably improve performance for abstract semantic similarity in small-scale archival settings.



% Section 5 Discussion
\section{Discussion}

\subsection{RQ1: Capacity Effects in Small-Scale Retrieval}
The results indicate that embedding-model capacity does not consistently improve retrieval performance in a small-scale classical music archive. 

This suggests that the benefits of scaling are highly task-dependent, particularly for abstract semantic tasks such as musical character retrieval. This observation is consistent with prior work, which shows that pretrained audio models exhibit substantial performance variability across MIR and recommendation tasks, and that strong MIR performance does not necessarily translate to improved recommendation quality \cite{tamm2024}.

Furthermore, classical music within a single archive exhibits strong stylistic coherence,
which may lead to reduced representational separability in the embedding space. We emphasize that this is a hypothesis rather than a directly validated property of the embeddings, as no explicit geometric analysis (e.g., pairwise distance statistics) is conducted in this study. Future work should investigate embedding-space structure more directly to validate this effect. This effect is further influenced by the relationship between the number of relevant items and the evaluation cutoff. In particular, when a large number of relevant items exist per query due to multi-label assignments, set-based metrics such as Recall and F1 become less informative at small cutoff values, as only a small fraction of the relevant set can be retrieved.

\subsection{RQ2: Conditions for Capacity Benefits}
The results suggest that higher-capacity models may provide limited benefits when the retrieval task is aligned with structured and well-defined attributes, such as composer identity. However, these effects are not consistent across model families.

In contrast, for abstract semantic tasks such as musical character retrieval,
the benefit of scaling diminishes. This may be due to the coarse and subjective nature of the labels, which limits the ability of larger models to learn more discriminative representations.

\subsection{Hypothesis Evaluation}
The results partially support the initial hypothesis that increasing embedding-model capacity does not necessarily improve retrieval performance. However, the findings suggest a more nuanced interpretation: capacity effects are strongly conditioned by the structure of the retrieval task.
While increased capacity provides measurable benefits for structured tasks such as composer-based retrieval, it does not consistently improve performance for abstract semantic tasks such as musical-character retrieval. Therefore, the hypothesis is refined to a condition-dependent form, where the effectiveness of model scaling is heavily governed by task definition, though future work is needed to determine how dataset characteristics mediate this relationship.

\subsection{Implications for Backend Design in Classical Music Archives}
From a system design perspective, the results suggest that, for small-scale classical music archives, increasing embedding-model capacity is not the most effective strategy for improving retrieval performance. Instead, lightweight or mid-sized models may provide comparable performance at lower computational cost. These findings indicate that system design may benefit from prioritizing task definition, label quality, and dataset curation over model scaling. In practical archival systems such as iPalpiti, this suggests that efficient models may be sufficient for real-time retrieval and deployment-constrained environments. While this study focuses on item-to-item similarity retrieval to isolate embedding quality, real-world recommender systems typically operate on user profiles rather than single query items. Therefore, future work should investigate how these representations perform when aggregated into user-level preference models in cold-start recommendation settings.

\subsection{Limitations}
This study has several limitations. First, the dataset is small and restricted to a single-domain classical music archive, which limits the generalizability of the findings. Additionally, the small size of the dataset (203 items) inherently limits the diversity of the candidate pool during retrieval, which may restrict the ability of high-capacity models to fully demonstrate their discriminative advantages. Second, the musical-character labels used in the primary proxy task are coarse and partially subjective, which constrains the evaluation of fine-grained semantic similarity. Third, the evaluation relies on offline proxy tasks rather than real user interaction, and therefore does not directly measure recommendation effectiveness in practical settings. Finally, the pseudo-labels used for the musical-character task are partially derived from a pretrained emotion recognition model. This introduces a potential circularity in the evaluation, as the same class of pretrained representations is used both to generate labels and to evaluate embedding quality. As a result, the observed performance may reflect alignment with the labeling model rather than true semantic similarity.

\subsection{Directions for Future Work}
Future work should investigate larger and more diverse music datasets to evaluate whether the observed capacity limitations persist at scale. In addition, incorporating richer and more fine-grained semantic labels may improve the evaluation of abstract musical similarity. Finally, integrating user interaction data and conducting user studies would provide a more direct assessment of retrieval effectiveness in real-world archival systems.

% Section 6 Conclusion
\section{Conclusion}
This study examined the effect of embedding-model capacity on content-based music similarity retrieval in a small-scale classical music archive. The results show that increasing model capacity does not consistently improve retrieval performance. While higher-capacity models may provide limited gains for structured tasks, these effects are not consistent across model families, and scaling effects are weak or absent for abstract semantic retrieval.

These findings indicate that retrieval performance in small-scale settings is strongly conditioned by the structure of the retrieval task rather than by embedding-model capacity alone. In particular, we show that evaluation outcomes are highly sensitive to the relationship between the number of relevant items and the evaluation cutoff. When a large set of relevant items exists per query, set-based metrics such as Recall and F1 provide limited signal at small cutoff values, while ranking-based metrics such as NDCG and Hit@K remain more stable.

Overall, this study highlights the importance of aligning model capacity, task definition, and evaluation design when developing content-based retrieval systems for data-constrained music archives.

% Generate references from the BibTeX database.
\bibliography{ISMIRtemplate}

\end{document}
